\section{Tapered iterative rounding}\label{sec:rounding}

This section rounds a sufficiently fine grid solution.  The graph used for
rounding gives every active copy one unit of possible incoming erasers.
Its distance-dependent taper makes the opening that accompanies an erasure
comparable to the client's surviving candidates.  A potential that retains
some cheap candidates after an opening converts this comparison into an
application of Theorem~\ref{thm:finite-inequality}.

\begin{theorem}[Rounding a grid solution]\label{thm:grid-rounding}
For every fixed $\tau>0$ and all sufficiently large constant integers
$\Delta$, put $g=1/\Delta$.  Suppose a finite rational metric instance has openings
$\widetilde y_i\in[0,1]\cap g\mathbb Z$ of total mass $M\ge1$.
For a client $j$, let $V_{j,0}$ be the cost of its nearest unit of this mass.
There is a constant $C=C(\tau,\Delta)$ and a randomized algorithm of
expected polynomial running time that outputs a set $O$ of facilities with
\[
 |O|\le M+C,
 \qquad
 \E d(j,O)\le(1+2/e+\tau)V_{j,0}\quad\text{for every client }j.
\]
All parameters and the additive constant are independent of the instance.
In particular, an initial mass at most $k+c_0$, for a fixed constant
$c_0$, gives at most $k+c_0+C$ facilities.
The deterministic implementation of the resulting total-cost guarantee
is given in Section~\ref{sec:derandomization}.
\end{theorem}

We prove the theorem using the functions $A,P,B,m$ defined in the preceding
section.  Thus $B(x)=2P(x)=A'(x)$ and
$m(x)=A(x)-1-xB(x)$.  We use, in particular,
\begin{equation}\label{eq:rnd-function-identities}
 \begin{gathered}
 0<B(x)\le1,\qquad B(x)\ge2P(1)>0,\qquad m(x)\ge0,\\
 (1+x)A(x)+xB(x)-1=3x\qquad(0\le x\le1).
 \end{gathered}
\end{equation}
The identities and derivative bounds used below follow directly from
the defining integrals.  Constants implicit in $O_\Delta(\cdot)$ may
depend on $\Delta$, but not on the instance, the round, or the subsequent
heavy-source threshold.

\subsection{The graph and the physical process}

The source-opening and neighbor-erasure operation follows the framework
of Byrka et al.~\cite[Definitions~3.2 and~5.1]{ByrkaEtAl2026}.
We specify the graph, taper, and batches here because their precise
properties will be used in the sub-two drift analysis.

Represent each opening by active copies of mass $g$, grouped by facility.
Each group has at most $\Delta$ active copies.  Co-located facilities
introduced by earlier splitting may remain distinct groups throughout
the process; they can be merged in the output. A group always means one
such labeled facility group, rather than an entire zero-distance class.
An integral group has exactly $\Delta$ copies.

Fix a small time step $\ell>0$ and a large time horizon $\Lambda$.
There are $T=\lceil\Lambda/\ell\rceil$ rounds.  At the start of round
$r\in\{0,\ldots,T-1\}$ let $\xi=r\ell$ and $x=e^{-\xi}$.
For that round choose a fixed rational number $\widehat P$ satisfying
\begin{equation}\label{eq:rnd-taper-approximation}
 0<\widehat P\le\tfrac12,
 \qquad |\widehat P-P(x)|\le\zeta,
 \qquad 0<\zeta<\tfrac14.
\end{equation}
The finitely many rational approximations may be fixed in advance.
The analytic value $x$ is used only in the analysis.

For each active target copy $i$, independently of the other targets
conditional on the current state, choose a set $\mathcal H_i$ of exactly
$\Delta$ active source copies.  If the colocated active mass is below one,
choose $N_i>0$ minimally so that
\begin{equation}\label{eq:rnd-edge-marginals}
 q_{hi}=\min\left\{1,\frac{(1-d(i,h)/N_i)_+}{\widehat P}\right\},
 \qquad \sum_h gq_{hi}=1,
\end{equation}
and give $\mathcal H_i$ these marginals.  Such a radius exists: the sum
is continuous and nondecreasing in the positive radius, starts below
one, and reaches the total active mass at a finite radius.  This last
observation also handles total active mass exactly one.  The equation
is piecewise linear in $1/N_i$.  Fixed-sum rounding realizes the
marginals since their unscaled sum is the integer $\Delta$.

If the colocated active mass is already at least one, set $N_i=0$ and
choose $\Delta$ colocated copies, including every active copy in the
target's own group.  Any fixed such choice suffices in this case; denote
its marginals by $q_{hi}$ as well.  In either case every own-group copy
belongs to $\mathcal H_i$ with certainty, and
\begin{equation}\label{eq:rnd-taper-comparison}
 \begin{aligned}
  N_i(1-\widehat Pq_{hi})&\le d(i,h)&&\text{if }q_{hi}<1,\\
  d(i,h)&\le N_i(1-\widehat Pq_{hi})&&\text{if }q_{hi}>0.
 \end{aligned}
\end{equation}
For $N_i>0$ these statements follow from the three cases of the clipped
linear formula; for $N_i=0$ all sources of positive marginal are
colocated.  Define the out-set of source $h$ by
\[
 E_h=\{i:h\in\mathcal H_i\}.
\]
A full update at $h$ erases $E_h$ and opens the facility of $h$.

We use balanced and biased unbalanced batches, with forced openings for
large updates, as in the approach of Byrka et
al.~\cite[Algorithms~4--5 and Section~5.4]{ByrkaEtAl2026}.
The exact batch and chunk rules below are part of our construction.
To implement many updates while controlling the total opening mass, put
\[
 u_h=1-g|E_h|,
 \qquad
 \mathcal Z=\{h:u_h=0\},\quad
 \mathcal P=\{h:u_h>0\},\quad
 \mathcal N=\{h:u_h<0\}.
\]
Every target has $\Delta$ incoming sources, so $\sum_hu_h=0$.  Write
\begin{equation}\label{eq:rnd-balance-size}
 A_* =\sum_{h\in\mathcal P}u_h
     =\sum_{h\in\mathcal N}|u_h|,
 \qquad
 |\mathcal P\cup\mathcal N|\le 2A_*/g.
\end{equation}
The last inequality uses that each nonzero $u_h$ is a nonzero multiple
of $g$.  Choose a further small constant $b_*>0$, to be fixed last, and put
\[
 B_* =\max\{1,b_*A_*\}.
\]
A source is \emph{heavy} if $g|E_h|>B_*$.  Force open the facilities of
all heavy sources, whether or not any random update is triggered.

All random choices and erasure sets in a round refer to its pre-round
graph and active copies.  Set $p_*=2\ell g$, and fix a small bias
$0<\phi\le1$.  Flip an independent fair coin to choose one of the following
two types of round.
\begin{enumerate}
\item In a balanced round, propose each $h\in\mathcal Z$ independently
with probability $p_*$.  Select a proposed source only if no other
proposed source has an out-set intersecting its out-set.
\item In an unbalanced round, select sources of $\mathcal P$
independently with probability $p_*$, and nonheavy sources of
$\mathcal N$ independently with probability $p_*(1+\phi)$.
For each heavy source, partition its out-set into chunks of mass at
most $B_*+1$ and trigger each chunk independently with probability
$p_*(1+\phi)$.  A triggered chunk erases just that chunk; its source
facility is already included among the forced openings.
\end{enumerate}
All probabilities are at most one for the parameter choices below.
Take the union of the selected full out-sets and triggered chunks and
erase it.  Then reset each forced or selected facility to $\Delta$
active copies.  Copies removed by either erasure or reset are distinct
from the regenerated copies.

Two invariants are useful.  First, an integral target group's own
$\Delta$ copies exhaust each of its incoming sets.  Therefore an
erasure of an integral group is caused by one of its own sources; that
source either is selected or is heavy.  The group is reset in either
case.  Thus integral facilities persist, including when other facilities
are colocated with them.  Second, every erasure has an associated
selected or forced opening.  Hence total active mass always stays at
least one.  The construction of the next round's graph is consequently
well defined.

\subsection{The opening-count bound}

We first bound the number of facilities independently of connection cost.
The next lemma controls the total active mass with high probability;
completion will later turn that mass bound into a facility-count bound.

\begin{lemma}[Constant additive count]\label{lem:rnd-count}
Fix $g,\phi,\ell,T$ as above, with $\ell$ sufficiently small as a
function of $\Delta$ and $\phi$.  For every fixed $\eta>0$, one can
choose $b_*>0$ and a constant $C$ such that the final active mass is at
most $M+C$ with probability at least $1-\eta$.  The bound holds
uniformly over all starting states of total mass $M$.
\end{lemma}

\begin{proof}
Condition on the entire pre-round state and its realized graph.
Write $A=A_*$, $B=B_*$, $p=p_*$, and $m_h=g|E_h|$ for this proof.
If $n$ is the number of active copies, then $\sum_hm_h=n$.
Each source in the zero class has $m_h=1$, so removing that class gives
\[
 \sum_{h\in\mathcal P\cup\mathcal N}m_h
 =|\mathcal P\cup\mathcal N|.
\]
Every heavy source lies in $\mathcal N$ because $B\ge1$.  If $H$ is
their number, then
\begin{equation}\label{eq:rnd-heavy-count}
 H\le\frac{|\mathcal P\cup\mathcal N|}{B}
 \le\frac{2A}{g\max\{1,b_*A\}}
 \le\frac{2}{gb_*}.
\end{equation}
Thus forced resets add at most a constant mass.  If $A=0$, there are
no heavy sources and the same conclusions hold directly.

In a balanced round, the selected erasure sets are disjoint and each
has mass one.  Opening one facility for each selected source adds at
most one unit per set, so there is no mass increase except the forced
allowance.

For an unbalanced round index the ordinary selections and heavy
chunks by actions $e$.  Let $I_e$ be the trigger indicator, $S_e$ its
erased set, $m_e=g|S_e|$, and $a_e=1$ for an ordinary selection or $0$
for a chunk.  Resetting facilities adds at most one unit per ordinary
selection in addition to the forced allowance. More explicitly, the
change in active mass is at most
\[
 H+\sum_e a_e I_e
   -g\left|\bigcup_{e:I_e=1}S_e\right|.
\]
The union is a set of pre-round copy identities. A reset may remove
additional old copies outside this union, which only decreases the
change. Counting a selected or forced group more than once likewise
only enlarges the upper bound. Truncated inclusion--exclusion bounds
the erased union from below. Consequently the mass increase, apart
from forcing, is at most
\begin{equation}\label{eq:rnd-count-statistic}
 Z_* =\sum_e I_e(a_e-m_e)
       +\sum_{e<f}g|S_e\cap S_f|I_eI_f.
\end{equation}
This remains an upper bound when several selected copies have the
same facility or when a selected facility was also forced.

Chunking preserves total erasure mass.  There are
$|\mathcal P\cup\mathcal N|-H$ ordinary actions, and hence
\begin{equation}\label{eq:rnd-action-total}
 \sum_e(m_e+a_e)
 =2|\mathcal P\cup\mathcal N|-H\le4A/g.
\end{equation}
The expectation of the linear part of \eqref{eq:rnd-count-statistic}
is exactly
\[
 pA-p(1+\phi)A-p(1+\phi)H\le-p\phi A.
\]
The last term pays for the absence of a new opening in a chunk action.
Each target belongs to at most $\Delta$ action sets, even after
chunking.  Therefore
\[
 \sum_{e<f}g|S_e\cap S_f|
 \le\frac{\Delta-1}{2}\sum_em_e.
\]
Independence of the indicators bounds the overlap expectation by
$O(p^2\Delta A/g)$.  Since $p=2\ell g$, taking $\ell$ small enough
relative to $\phi$ and $\Delta$ gives
\begin{equation}\label{eq:rnd-count-mean}
 \E Z_*\le-p\phi A/2.
\end{equation}

Flipping coordinate $I_e$ changes $Z_*$ in absolute value by at most
$a_e+\Delta m_e$.  The Efron--Stein independent-coordinate variance inequality
\cite[Section~2.2, Proposition~1]{BBLM2005}, the bound
$a_e+m_e\le B+2$, and \eqref{eq:rnd-action-total} give
\begin{equation}\label{eq:rnd-count-variance}
 \Var Z_*
 \le O(p\Delta^2)\sum_e(a_e+m_e)^2
 \le O(p\Delta^2)(B+2)A/g.
\end{equation}
If $A>1/b_*$, then $B=b_*A>1$.  Combining
\eqref{eq:rnd-count-mean} and \eqref{eq:rnd-count-variance} yields
\begin{equation}\label{eq:rnd-count-failure}
 \Prob(Z_*>0)
 \le O\!\left(\frac{\Delta^2b_*}{p\phi^2g}\right).
\end{equation}
If instead $A\le1/b_*$, there are at most $2/(gb_*)$ ordinary
sources in the nonzero classes.  The actual positive mass change is
then bounded deterministically by this constant plus the forced
allowance; a deterministic bound on the quadratic statistic is
unnecessary.

Choose $b_*$ sufficiently small that the union bound for
\eqref{eq:rnd-count-failure} over the $T$ rounds is at most $\eta$.
On the complementary event the total increase is at most, for example,
$4T/(gb_*)$.  This proves the lemma with a constant $C$.
The proof was conditional on every state and graph, so the same union
bound applies to the adaptive sequence of rounds.
\end{proof}

For later use, the mean and variance of $Z_*$ depend only on indicator
moments through degree four, since $Z_*$ is a quadratic polynomial.
The probability bounds for balanced selections below need only
first and second moments of their proposal indicators.  These are
the moment requirements used in Section~\ref{sec:derandomization}.

\subsection{A persistent backup and final completion}

We next establish a persistent bound on each client's eventual
connection distance, including trajectories on which none of its assigned
copies opens. This is the backup-and-completion role in iterative
rounding~\cite[Section~3.3 and Lemmas~6.3 and~6.5]{ByrkaEtAl2026};
the following argument proves the bound for our tapered graph and reset
rule.

Fix a client for the remainder of the analysis.  Let $S_0$ be the
$\Delta$ copies in its initial nearest unit, chosen using any fixed
tie-breaking rule.  Write $d_i$ for its distance to copy $i$,
\[
 V_0=\sum_{i\in S_0}gd_i,
 \qquad R=\max_{i\in S_0}d_i.
\]
In particular $R\le\Delta V_0$.

\begin{lemma}[Backup and completion]\label{lem:rnd-completion}
Once a copy of $S_0$ is physically removed, an integral facility within
distance $5R$ of the client exists and persists to the end.
At the end of the process one can retain all integral facilities and
open further facilities, using no more than the total final active
mass, so that every client is connected within its own distance $5R$.
\end{lemma}

\begin{proof}
Consider the round of the first physical removal from $S_0$.  Just
before that round all of $S_0$ is still active.  For $i,t\in S_0$,
$d(i,t)\le d_i+d_t\le2R$.  At radius $2R/(1-\widehat P)$, the
$\Delta$ copies of $S_0$ all have incoming marginal one for target
$i$.  Minimality of the taper radius gives
\[
 N_i\le2R/(1-\widehat P)\le4R.
\]
If $R=0$, the zero-radius rule gives the same conclusion directly.
An erasing source $h$ has $d(i,h)\le N_i$ by
\eqref{eq:rnd-taper-comparison}, and hence $d_h\le5R$.
Its facility is either selected or forced open.  A removal by reset
of $i$'s own group instead opens a facility within $R$.
The persistence invariant proves the first assertion.

At the end, consider the clients not already served within $5R$ by
an integral facility.  The first assertion implies that every such
client still has all its original $S_0$ copies.  Its radius-$R$
neighborhood of active copies therefore has mass at least one.
That neighborhood contains no integral group, since such a group
would itself provide the required connection.

Process these neighborhoods in nondecreasing order of their radii.
Accept a neighborhood if it is disjoint from all previously accepted
ones and open any facility in each accepted neighborhood.  These
disjoint neighborhoods each contain at least one unit of mass and
contain no integral group.  Their openings, together with the retained
integral groups, therefore fit within the final active mass.
If a radius-$R$ neighborhood was skipped, it meets an accepted
radius-$R'$ neighborhood with $R'\le R$.  Through an intersection
point, the distance to the accepted neighborhood's chosen facility
is at most $R+2R'\le3R$.  This proves the asserted $5R$ bound for
every client, also when some radii are zero.
\end{proof}

At any intermediate time define $b$ to be the distance from the fixed
client to its nearest integral facility, capped at $5R$; when no
integral facility exists use $b=5R$.  By persistence and
Lemma~\ref{lem:rnd-completion}, the eventual completed solution has
connection distance at most $b$ at every time on the trajectory.
The cap need not already be attained by an open facility.  It is a
valid bound on the eventual connection because completion supplies
$5R$ on every trajectory.  The quantity $b$ never increases.

\subsection{The potential and its single-action bounds}

For analysis only, track an alive set $S\subseteq S_0$ of original
copies.  A physically removed copy leaves $S$ permanently, even when
its facility is reset with regenerated copies.  We also permit the
additional discards specified below.  Put
\[
 y=g|S|,\qquad V=\sum_{i\in S}gd_i.
\]
At elapsed time $\xi$, with $x=e^{-\xi}$, define
\begin{equation}\label{eq:rnd-potential}
 f=aV+(1-y)b,
 \qquad
 a=A(x)+B(x)(y-x)=1+m(x)+B(x)y.
\end{equation}
The coefficient $a$ is the value at $y$ of the tangent line to the
concave function $A$ at $x$; $x$ is a deterministic clock, while $y$
is the realized alive mass. Its affine dependence on $y$ makes the
erasure correction quadratic in removed mass and cost, so the graph
average of the ideal bound below requires only first and second edge
moments.
In particular $a\ge1$.  The initial potential is
\begin{equation}\label{eq:rnd-initial-potential}
 f_0=A(1)V_0=(1+2/e)V_0.
\end{equation}
If $C_j$ denotes the actual connection distance in the completed
output, the preceding cap property gives
\begin{equation}\label{eq:rnd-terminal-correction}
 C_j-f_T\le y_Tb_T-a_TV_T\le5Ry_T.
\end{equation}
The potential and every single-action bound below are $O(R)$ uniformly
over time: $V\le R$, $0\le y\le1$, $b\le5R$, and $a$ is bounded.

We first freeze time and the graph.  Write $B=B(x)$ and
$b_h=\min\{b,d_h\}$.  For a source $h\notin S$, let
\[
 \delta_h=g|E_h\cap S|,
 \qquad Y_h=\sum_{i\in E_h\cap S}gd_i.
\]
If only this full source update affects $S$, its removal changes the
coefficient to $a-B\delta_h$ and makes the new cap at most $b_h$.
Its potential change is consequently at most
\begin{align*}
 &(a-B\delta_h)(V-Y_h)+(1-y+\delta_h)b_h-f\\
 &\quad=\delta_hb_h-aY_h-B\delta_h(V-Y_h)
          +(1-y)(b_h-b)\\
 &\quad\le\delta_hb_h-aY_h-B\delta_h(V-Y_h).
\end{align*}
The same statement holds when $h$ is not alive but its group contains
alive copies: all those copies belong to $E_h$ by the own-group edges,
so its reset creates no unaccounted removal.

If an alive source $h\in S$ opens, retain among the physical survivors
only copies satisfying $ad_i<d_h$.  If their mass and cost are $y',V'$,
their new coefficient is at most $a$ and their new cap is at most
$d_h$.  Therefore
\[
 f'\le aV'+(1-y')d_h
 =d_h-\sum_{\substack{i\in S\\h\notin\mathcal H_i}}
            g(d_h-ad_i)_+.
\]
Copies satisfying equality can be discarded without changing this
bound. Figure~\ref{fig:survivors} separates physical erasure from these
additional analysis discards. This continuation after an opening
supplies the subtraction that will enter
Theorem~\ref{thm:finite-inequality}.

\begin{figure}[tbp]
\centering
\begin{tikzpicture}[
  box/.style={draw,rounded corners=2pt,align=center,inner sep=7pt,
              text width=3.95cm,font=\small},
  flow/.style={-{Stealth[length=2mm]},thick},
  note/.style={align=center,font=\small}]
\node[note] (source) {A full update at an alive source $h\in S$ opens its facility};
\node[box,below=10mm of source,xshift=-4.8cm] (erased)
  {$S\cap E_h$\\[3pt]Physically erased\\and removed from $S$};
\node[box,right=4mm of erased] (kept)
  {$i\in S\setminus E_h$\\with $ad_i<d_h$\\[3pt]Physically survives\\and remains alive};
\node[box,right=4mm of kept] (discarded)
  {$i\in S\setminus E_h$\\with $ad_i\ge d_h$\\[3pt]Physically survives;\\discarded from the analysis};
\draw[flow] (source.south) -- ++(0,-4mm) -| (erased.north);
\draw[flow] (source.south) -- (kept.north);
\draw[flow,dashed] (source.south) -- ++(0,-4mm) -| (discarded.north);
\node[box,text width=12.95cm,below=8mm of kept] (reset)
  {The opened facility is reset to $\Delta$ active copies.\\
   These regenerated copies are not original identities in $S_0$ and do not enter $S$.};
\end{tikzpicture}
\caption{One full source update, with no simultaneous relevant action.
Only the middle group contributes to the continued alive assignment.
The dashed branch is an analysis discard, not a physical erasure.
Positions are schematic and represent no metric geometry.}
\label{fig:survivors}
\end{figure}


The ideal first-order bound, assigning rate $g$ to each full source
update, is thus
\begin{equation}\label{eq:rnd-ideal-generator}
\begin{aligned}
 \mathcal B={}&
 \sum_{h\in S}g\left[
 d_h-\sum_{\substack{i\in S\\h\notin\mathcal H_i}}
                  g(d_h-ad_i)_+-f\right]\\
 &+\sum_{h\notin S}g\left[
 \delta_hb_h-aY_h-B\delta_h(V-Y_h)\right].
\end{aligned}
\end{equation}
At a fixed alive set and cap, the time derivative is
\begin{equation}\label{eq:rnd-time-derivative}
 \dot f=-x(y-x)B'(x)V.
\end{equation}
Indeed the derivative of $A(x)+B(x)(y-x)$ with respect to $x$ is
$B'(x)(y-x)$.

\subsection{From the physical batch to the ideal bound}

For the actual analysis update, first account for forced openings by
removing from $S$ all copies whose groups are forced and updating the
cap.  Then apply the round's erasures and resets.  If a source that
was alive after force removals but before these erasures was selected,
choose one such source $h$ and discard every remaining copy with
$ad_i\ge d_h$, using the pre-round value of $a$.  This is purely an
analysis rule and need not be computed by the algorithm.

\begin{lemma}[Conditional batch drift]\label{lem:rnd-batch}
For every pre-round history and every realized graph, the physical
round, the above alive-set update, and the time advance satisfy
\begin{equation}\label{eq:rnd-batch-drift}
 \E[f_{r+1}-f_r\mid\text{history, graph}]
 \le\ell(\dot f+\mathcal B)
       +O_\Delta\bigl((\phi+\ell)\ell R\bigr).
\end{equation}
Here $\mathcal B$ is computed from that graph and the original
pre-round alive set.  The constants and the required upper bound on
$\ell$ are uniform in $b_*$.
\end{lemma}

\begin{proof}
First suppose no group containing an alive copy is forced.  Call an
action relevant if its ordinary source is in $S$, or its ordinary
out-set or heavy chunk meets $S$.  There are at most
$\Delta^2+\Delta$ relevant actions: each of at most $\Delta$ alive
targets has $\Delta$ incoming sources.  Own-group edges ensure that
this list also includes every ordinary reset that can affect $S$.

A balanced source has an out-set of exactly $\Delta$ copies.  Every
one of those copies belongs to at most $\Delta$ source out-sets, so
the number $c_h$ of potentially conflicting balanced proposals is at
most $\Delta^2$.  Including the type coin, its selection marginal is
\[
 \tfrac12p_*(1-p_*)^{c_h}
 =\ell g\bigl(1+O_\Delta(\ell)\bigr).
\]
An unbalanced ordinary or chunk action has marginal $\ell g$ or
$(1+\phi)\ell g$.  Any two distinct actions have joint probability
$O(\ell^2g^2)$; actions of opposite types have joint probability zero.
These estimates do not assert independence of balanced selections
after conflict filtering.

If no relevant action occurs, only the cap can improve, so the
frozen-time potential cannot increase.  If exactly one relevant
ordinary action occurs, its change is bounded by the appropriate
single-source expression in \eqref{eq:rnd-ideal-generator}.
Unrelated selections or forced openings can only lower its cap and
do not invalidate that upper bound.  The probability of two or more
relevant actions is $O_\Delta(\ell^2)$, and all possible potentials
and individual bounds have magnitude $O(R)$.  Replacing isolated-action
probabilities by marginals and their marginals by $\ell g$ therefore
costs $O_\Delta((\phi+\ell)\ell R)$.

For a heavy source outside $S$, its forced facility makes the cap at
most $b_h$ even when only one of its chunks triggers.  Force accounting
can be done before erasure accounting because no forced group in the
present case meets $S$; it only improves the cap.  Write
$\delta_c,Y_c$ for the alive mass and cost removed by chunk $c$ of
this source, and $\delta=\sum_c\delta_c$, $Y=\sum_cY_c$.  Its summed
single-chunk bounds satisfy
\begin{align*}
 \sum_c[\delta_cb_h-aY_c-B\delta_c(V-Y_c)]
 &=\delta b_h-aY-B\delta V+B\sum_c\delta_cY_c\\
 &\le\delta b_h-aY-B\delta(V-Y).
\end{align*}
All cross terms omitted in the last step are nonnegative.  Consequently
the chunk bounds are dominated by the full-source term used in
$\mathcal B$.

The functions in the potential have uniformly bounded derivatives.
The chance that a relevant action changes $S$ is $O_\Delta(\ell)$;
changes of the cap alone do not affect \eqref{eq:rnd-time-derivative}.
Advancing time thus adds at most
$\ell\dot f+O_\Delta(\ell^2R)$.  This proves the lemma when no alive
group is forced.

Now suppose that forcing does remove alive copies.  Let their total
mass and cost be $\delta\ge g$ and $Y$, and let $b'\le b$ be the
cap after force accounting.  Every removed copy is in an opened
facility, so $Y\ge\delta b'$.  At frozen time, the resulting decrease
is at least
\begin{align}
 f-f'
 &\ge(a-1)Y+B\delta(V-Y)+(1-y)(b-b')\notag\\
 &\ge BgV+(1-y)(b-b').\label{eq:rnd-force-savings}
\end{align}
For the second line, use $a-1=m+By$, $m\ge0$, and
$y\ge\delta\ge g$.  At least one removed alive copy has distance
at most $V/g$, so $b'\le V/g$.  With
\[
 Q=V+(1-y)b,\qquad D=(1-y)(b-b'),
\]
we obtain
\[
 \begin{gathered}
 Q\le(1+1/g)V+D,\qquad BgV+D\ge c_\Delta Q,\\
 c_\Delta=\min\left\{1,\frac{2P(1)g}{1+1/g}\right\}>0.
 \end{gathered}
\]
Thus forcing supplies a fixed fraction of $Q$ as free savings.

After force accounting, the remaining alive cost is at most $V$ and
the cap is at most $V/g$.  Any further frozen-time increase has
magnitude $O_\Delta(V)$, and a change to the alive set still has
probability $O_\Delta(\ell)$.  If the alive set does not change, the
cap only improves.  Subsequent actions and time advance therefore
have expected increase at most $O_\Delta(\ell Q)$.

For comparison with the pre-force ideal expression, we also have
\begin{equation}\label{eq:rnd-force-generator-size}
 |\dot f|+|\mathcal B|\le O_\Delta(Q).
\end{equation}
For inside sources, $d_h\le V/g$.  If $y<1$, the grid gives
$1-y\ge g$, so an external cap charge is bounded using $b\le Q/g$.
If $y=1$, the alive set itself is a full unit; its maximum distance
is at most $V/g$.  The radius argument from
Lemma~\ref{lem:rnd-completion} then bounds any source erasing an alive
copy within $5V/g$ of the client.  Only $O_\Delta(1)$ sources have
nonzero contributions.  These observations prove
\eqref{eq:rnd-force-generator-size}.

Choose $\ell$ sufficiently small in terms of these constants and
$c_\Delta$.  The savings in \eqref{eq:rnd-force-savings} absorb both
the subsequent expected increase and the possible negative quantity
$\ell(\dot f+\mathcal B)$.  This proves
\eqref{eq:rnd-batch-drift} also in the forced case, without an
additional error term.  If $Q=0$, all the quantities in this
comparison vanish, so no division by $Q$ is needed.  More explicitly,
if $V=0$ then $b'=0$, and the force savings remove $(1-y)b$; when
also $y=1$, relevant source distances and all generator terms are
zero.  The argument is uniform in the later choice of $b_*$.
\end{proof}

The order of conditioning in Lemma~\ref{lem:rnd-batch} is essential.
Whether a group is forced depends on the graph.  We have proved the
drift comparison for \emph{every} realized graph, in both force cases;
we will now average that comparison over the original graph law,
without conditioning the law on a forcing event.

\subsection{Averaging the graph and applying the finite inequality}

The physical batch has now been compared with the ideal generator for
every realized graph. We average that generator using edge marginals and
cross-target second moments, then verify the hypotheses of the abstract
finite comparison inequality. The resulting bound is the cost estimate
that will be summed over rounds.

\begin{lemma}[Averaged ideal drift]\label{lem:rnd-average}
Conditional on the pre-round history, the ideal expression satisfies
\begin{equation}\label{eq:rnd-average-drift}
 \dot f+\E_{\mathrm{graph}}\mathcal B
 \le BgV+12\zeta V.
\end{equation}
\end{lemma}

\begin{proof}
The pre-round quantities $S,y,V,b,x$ are fixed throughout this proof.
Define
\begin{equation}\label{eq:rnd-graph-statistics}
 \begin{aligned}
 w_i&=\sum_{t\in S}gq_{ti},\\
 K_{it}&=\sum_{h\notin S}gq_{hi}q_{ht},\\
 s_i&=\frac{\sum_{h\notin S}gq_{hi}^2}{1-w_i}
       \quad\text{if }w_i<1,
       \qquad s_i=0\quad\text{if }w_i=1.
 \end{aligned}
\end{equation}
We have $0\le w_i\le y\le1$ and $0\le s_i\le1$.
The inequality $2uv\le u^2+v^2$ gives
\begin{equation}\label{eq:rnd-kernel-bound}
 K_{it}\le\tfrac12\bigl[s_i(1-w_i)+s_t(1-w_t)\bigr].
\end{equation}
No symmetry of the directed marginals $q_{hi}$ is assumed.

To average the positive quadratic erasure term, put
$X_{hi}=\ind_{\{h\in\mathcal H_i\}}$.  Choices for distinct targets
are independent.  Expanding the two removed quantities consequently
gives the exact formula
\begin{align}
 \E_{\mathrm{graph}}\sum_{h\notin S}g\delta_hY_h
 &=\sum_{i,t\in S}g^2K_{it}d_t
   +\sum_{i\in S}g^2d_i\sum_{h\notin S}gq_{hi}(1-q_{hi})\notag\\
 &\le\sum_{i,t\in S}g^2K_{it}d_t+gV.
 \label{eq:rnd-diagonal-correction}
\end{align}
Here $\sum_hgq_{hi}=1$ bounds the diagonal correction.

For the cap contribution, the external incoming mass at $i$ is
$1-w_i=(1-y)+(y-w_i)$.  Split it in these proportions.  On its first
part use $b_h\le b$, and on its second part use, whenever $q_{hi}>0$,
\[
 b_h\le d_h\le d_i+d(i,h)
       \le d_i+(1-\widehat Pq_{hi})N_i.
\]
This yields
\begin{equation}\label{eq:rnd-backup-average}
 \sum_{h\notin S}gq_{hi}b_h
 \le(1-y)b+(y-w_i)\bigl[d_i+(1-\widehat Ps_i)N_i\bigr].
\end{equation}
If $w_i=1$, then $y=1$ and both sides are zero, so the convention
for $s_i$ requires no division in that case.

For completeness, averaging \eqref{eq:rnd-ideal-generator} before
using the cap bound gives, apart from the additive term $BgV$,
\begin{align*}
 &[1-(1+y)a-By]V-y(1-y)b\\
 &\quad+\sum_{i\in S}g\sum_{h\notin S}gq_{hi}b_h
       +a\sum_{i\in S}gw_id_i+BV\sum_{i\in S}gw_i
       +B\sum_{i,t\in S}g^2K_{it}d_t\\
 &\quad-\sum_{i,t\in S}g^2(1-q_{ti})(d_t-ad_i)_+.
\end{align*}
The first term of \eqref{eq:rnd-backup-average}, summed over $i$,
cancels $-y(1-y)b$.  The coefficient of $V$ after adding
\eqref{eq:rnd-time-derivative} simplifies using
\begin{equation}\label{eq:rnd-coefficient-identity}
 (1+y)a+yB+x(y-x)B'-1=3y+B(y-x)^2.
\end{equation}
Indeed \eqref{eq:rnd-function-identities} and its derivative give
\[
 (1+x)A+xB-1=3x,
 \qquad A+(2+x)B+xB'=3.
\]
Expanding the left side of \eqref{eq:rnd-coefficient-identity} in
$y-x$ has these as its constant and linear coefficients, and has
quadratic coefficient $B$.

It follows that $\dot f+\E_{\mathrm{graph}}\mathcal B-BgV$ is at most
\begin{equation}\label{eq:rnd-expanded-drift}
\begin{aligned}
 &-[3y+B(y-x)^2]V\\
 &\quad+\sum_{i\in S}g\Bigg\{
 (y-w_i)\bigl[d_i+(1-\widehat Ps_i)N_i\bigr]+aw_id_i
       +BVw_i+B\sum_{t\in S}gK_{it}d_t\\
 &\hspace{43mm}
       -\sum_{t\in S}g(1-q_{ti})(d_t-ad_i)_+\Bigg\}.
\end{aligned}
\end{equation}

If $y=0$, the ideal generator and time derivative are zero.  If
$y>0$ but $V=0$, then every $d_i$ with $i\in S$ is zero.  For any
$i\in S$ with $N_i>0$, all alive sources are colocated with $i$, so
$q_{ti}=1$ for every $t\in S$ and $w_i=y$.  If $N_i=0$, its radius
term already vanishes.  Every term of
\eqref{eq:rnd-expanded-drift} is then zero.  Thus the lemma holds
in these degenerate cases without normalization.

Assume now $V>0$, and set
\begin{equation}\label{eq:rnd-normalization}
 \mu=V/y,\qquad
 D_i=d_i/\mu,\qquad
 r_i=w_i/y,\qquad
 R_i=(1-2\zeta)N_i/\mu,\qquad C_i=1-Ps_i.
\end{equation}
Give each index $i\in S$ probability $g/y$.  Expectations over these
indices, written $\E_i$ and $\E_t$, satisfy
$\E_iD_i=1$ and $r_i=\E_tq_{ti}$.
For $q=q_{ti}<1$, the taper approximation implies
\[
 (1-2\zeta)(1-Pq)\le1-\widehat Pq,
\]
because the difference is at least
$2\zeta(1-Pq)-\zeta q\ge2\zeta(1-q)\ge0$.
The triangle inequality and \eqref{eq:rnd-taper-comparison} give
\begin{equation}\label{eq:rnd-normalized-comparison}
 R_i(1-Pq_{ti})\le D_i+D_t\qquad(q_{ti}<1).
\end{equation}
These are precisely the radius hypotheses of
Theorem~\ref{thm:finite-inequality}.

The error in changing the taper coefficient is controlled without
any bound on individual normalized radii.  Multiply the first
inequality in \eqref{eq:rnd-taper-comparison} by
$(g/y)(1-q_{ti})$, sum over $t\in S$, and use
$1-\widehat Pq_{ti}\ge1/2$ and $d(i,t)\le d_i+d_t$.  This proves
\begin{equation}\label{eq:rnd-radius-error-control}
 (1-r_i)N_i/\mu\le2(D_i+1).
\end{equation}
Since
\[
 0\le(1-\widehat Ps_i)-(1-2\zeta)(1-Ps_i)\le3\zeta,
\]
the replacement of
$(1-r_i)(1-\widehat Ps_i)N_i/\mu$ by
$C_i(1-r_i)R_i$ costs at most $6\zeta(D_i+1)$, whose expectation
is at most $12\zeta$.

Divide \eqref{eq:rnd-expanded-drift} by $yV=y^2\mu$ and put
$\Theta=P(y-x)^2/y$.  By \eqref{eq:rnd-kernel-bound}, the normalized
kernel term is at most
\[
 B\E_{i,t}K_{it}D_t
 \le P\E_i\bigl[s_i(1-yr_i)(1+D_i)\bigr].
\]
After the radius replacement, the entire normalized expression is
therefore at most
\begin{align*}
 -3-2\Theta+\E_i\Big[&
 (1-r_i)(D_i+C_iR_i)+ar_iD_i+2Pyr_i\\
 &+Ps_i(1-yr_i)(1+D_i)
 -\E_t(1-q_{ti})(D_t-aD_i)_+\Big]+12\zeta.
\end{align*}
Using $a=1+m+2Py$ and $\E_iD_i=1$, this is exactly
\begin{align*}
 \E_i\Big[&C_i(1-r_i)R_i
 +r_i\{mD_i+(2-s_i)Py(D_i+1)\}-C_i(D_i+1)\\
 &-\E_t(1-q_{ti})(D_t-(1+m+2Py)D_i)_+\Big]
 -2\Theta+12\zeta.
\end{align*}
Theorem~\ref{thm:finite-inequality} bounds the expectation by
$2\Theta$.  Multiplication by $yV$, with $y\le1$, proves
\[
 \dot f+\E_{\mathrm{graph}}\mathcal B
 \le BgV+12\zeta yV\le BgV+12\zeta V,
\]
as claimed.
\end{proof}

\subsection{Survival, accumulated error, and parameter choices}

\begin{lemma}[Decay of the alive mass]\label{lem:rnd-survival}
For sufficiently small $\ell$ depending only on $\Delta$, every
client's analysis process satisfies
\begin{equation}\label{eq:rnd-survival}
 \E y_T\le(1-\ell/2)^T\le e^{-\Lambda/2}.
\end{equation}
\end{lemma}

\begin{proof}
Condition on a history, a graph, and a currently alive target.
If its group is forced, the original target identity is removed
certainly.  Otherwise its $\Delta$ incoming source copies correspond
to $\Delta$ distinct possible actions: each ordinary source gives
its selection action, and each heavy source gives the unique chunk
containing this target.  The marginal and pair estimates in the
proof of Lemma~\ref{lem:rnd-batch} give an erasure probability at least
\[
 \Delta\ell g\bigl(1-O_\Delta(\ell)\bigr)
       -O(\Delta^2\ell^2g^2)\ge\ell/2.
\]
Here the lower bound is the union bound with pair intersections
subtracted.  A simultaneous reset does not restore the original
copy identity, and additional analysis discards only reduce survival.
The conditional expected surviving mass is thus at most
$(1-\ell/2)y$.  Iteration, starting from $y_0=1$, and $T\ell\ge\Lambda$
prove the lemma.
\end{proof}

We finish the proof of Theorem~\ref{thm:grid-rounding}.
First run the process without rejecting count failures, and complete
every output by Lemma~\ref{lem:rnd-completion}.  Average
Lemma~\ref{lem:rnd-batch} over the graph and use
Lemma~\ref{lem:rnd-average}.  Since $V\le V_0$, $B\le1$, and
$R\le\Delta V_0$, summing over $T$ rounds gives
\[
 \E f_T\le
 \left[1+2/e+(\Lambda+1)
       \{g+12\zeta+C_\Delta\Delta(\phi+\ell)\}\right]V_0
\]
for a constant $C_\Delta$, provided $\ell\le1$.
Equations~\eqref{eq:rnd-terminal-correction} and
\eqref{eq:rnd-survival} consequently imply the unconditional estimate
\begin{equation}\label{eq:rnd-total-error}
 \E C_j\le
 \left[1+2/e+5\Delta e^{-\Lambda/2}
       +(\Lambda+1)\{g+12\zeta+C_\Delta\Delta(\phi+\ell)\}\right]V_0.
\end{equation}
All these estimates also apply when $V_0=0$: then $R=0$, the completed
connection cost is zero, and no normalization by $V_0$ is used.

Here is an explicit order of choices that avoids circular dependence.
It suffices to consider $0<\tau\le1$.
\begin{enumerate}
\item Choose a sufficiently large constant integer $\Delta$ and put
$g=1/\Delta$.  For each such choice take, for example,
\[
 \Lambda=\left\lceil2\log(40\Delta/\tau)\right\rceil+1.
\]
Then $5\Delta e^{-\Lambda/2}\le\tau/8$.  Taking $\Delta$ sufficiently
large also makes $(\Lambda+1)g\le\tau/8$, since
$g\log(1/g)\longrightarrow0$.  Any additional upper bound on $g$
required by the preparation stage may be imposed at the same time.
\item With $\Delta,\Lambda$ fixed, choose a small positive rational
$\zeta$ so that $12(\Lambda+1)\zeta\le\tau/8$.
Next choose a small positive rational $\phi$, and then a small positive
rational $\ell$, so that
\[
 C_\Delta\Delta(\Lambda+1)(\phi+\ell)\le\tau/8.
\]
Require in addition every upper bound on $\ell$ used in the count,
batch, and survival lemmas, and $p_*(1+\phi)\le1$.
The batch constants are uniform in the remaining threshold $b_*$,
so these requirements can be imposed before choosing that threshold.
This fixes the constant number of rounds $T=\lceil\Lambda/\ell\rceil$.
\item Put $\alpha=1+2/e$ and choose a fixed
$0<\eta\le\tau/(4(\alpha+\tau))$.
Finally choose rational $b_*>0$ as in Lemma~\ref{lem:rnd-count} so
that the final mass is at most $M+C$ with probability at least
$1-\eta$.
\end{enumerate}
The unconditional cost in \eqref{eq:rnd-total-error} is then at most
$(\alpha+\tau/2)V_0$ for every client.  Repeat the entire process on
failure of the final mass test.  The accepted output is the
unconditional output conditioned on that test, whose success
probability is at least $1-\eta$.  Since connection costs are
nonnegative,
\[
 \E[C_j\mid\text{mass test succeeds}]
 \le\frac{\alpha+\tau/2}{1-\eta}V_0
 \le(\alpha+\tau)V_0.
\]
Completion uses at most the final active mass, so an accepted output
has at most $M+C$ facilities.  This proves the count and cost assertions
of Theorem~\ref{thm:grid-rounding}.

For fixed accuracy all parameters are instance-independent constants.
The number of current facilities never exceeds the number in the
split input, and every group has at most $\Delta$ active copies.
The graph, its marginals, a fixed-sum rounding, the action sets,
their intersections, and the greedy completion can therefore be
computed in polynomial time.  The radius equations and all physical
probabilities are rational; the approximations in
\eqref{eq:rnd-taper-approximation} are fixed constants.  Only the
analysis uses the transcendental functions or alive-set discards.
The expected number of repetitions is at most $1/(1-\eta)$, giving
expected polynomial running time.  Section~\ref{sec:derandomization}
replaces the local random choices by polynomial-size lists and obtains
a deterministic output with the corresponding total-cost bound.
